% !TEX root = ../main.tex
\chapter{Projetos}
\label{cap:projetos}

Tudo no SAPHO acontece dentro de um \textbf{projeto}: uma pasta no disco governada por um arquivo \file{.spf} (\emph{SAPHO Project File}). O projeto reúne um ou mais processadores, os arquivos Verilog de integração (\emph{top-level} e \emph{testbench}) e os metadados que a AURORA usa para orquestrar compilação e simulação.

\section{Criando um projeto}
\label{sec:criando-projeto}

Clique em \botao{Novo Projeto} (na toolbar ou na tela de boas-vindas). O modal \textbf{Criar Novo Projeto} pede:

\begin{longtable}{@{}lp{0.64\linewidth}@{}}
\toprule
\textbf{Campo} & \textbf{Regras} \\
\midrule
\endhead
\textbf{Nome do Projeto} & Apenas letras, números, sublinhado (\code{\_}) e hífen (\code{-}). Sem espaços, acentos ou símbolos. \\
\textbf{Local do Projeto} & Somente leitura; preencha com o botão \botao{Procurar}, que abre o seletor de pastas do Windows. A AURORA lembra o último local usado. \\
\bottomrule
\end{longtable}

Clique em \botao{Gerar Projeto}. Se algum campo violar as regras, o diálogo \textbf{Entrada Inválida} explica o problema (\enquote{Nome do Projeto e Local não podem conter espaços ou símbolos especiais\ldots}).

\figurapendente{Modal Criar Novo Projeto preenchido}{Com o nome \code{MeuFiltro} e um local qualquer, antes de clicar em Gerar Projeto.}

No disco, é criada a estrutura \file{<Local>/<Nome>/<Nome>.spf}. Um projeto recém-criado ainda não tem processadores --- a árvore abre praticamente vazia, e o próximo passo natural é o \botao{Hub de Processadores} (Capítulo~\ref{cap:criando-processadores}).

Para o exemplo condutor deste manual, crie o projeto \textbf{MeuFiltro}.

\section{O arquivo \texttt{.spf}}
\label{sec:spf}

O \file{.spf} é um arquivo JSON legível, com dois blocos:

\begin{itemize}
  \item \code{metadata} --- nome do projeto, datas de criação/modificação/última abertura, nome do computador e versão da AURORA que o criou;
  \item \code{structure} --- a lista de \code{processors}, as pastas do projeto, o \code{topLevelFile}, o \code{testbenchFile} e as listas de arquivos Verilog sintetizáveis e de \emph{testbench}.
\end{itemize}

Você raramente precisará editá-lo à mão: a AURORA é a única escritora do arquivo e o mantém em dia a cada ação na interface. Os caminhos são gravados relativos à pasta do projeto, então o projeto inteiro pode ser copiado, movido de máquina ou versionado no Git sem quebrar --- ao abrir, a AURORA realinha os caminhos ao novo local.

\section{Abrindo um projeto}

Três caminhos:
\begin{itemize}
  \item \botao{Abrir Projeto} na toolbar ou na tela de boas-vindas, selecionando o \file{.spf};
  \item dois cliques no \file{.spf} pelo Explorer do Windows (a extensão é associada na instalação);
  \item a lista de \textbf{Recentes} na tela de boas-vindas (até dez projetos, também disponíveis na \emph{jumplist} do ícone da AURORA na barra de tarefas do Windows).
\end{itemize}

Ao abrir, a AURORA carrega a árvore, restaura a lista de processadores e registra a data de abertura. Se o \file{.spf} tiver sumido do caminho registrado, o diálogo \textbf{Projeto Não Encontrado} avisa e o item sai da lista de recentes.

\section{Arquivos do projeto: quem é quem}
\label{sec:arquivos-do-projeto}

Um projeto maduro tem esta cara (exemplo com o nosso \code{MeuFiltro}):

\begin{lstlisting}
MeuFiltro/
|-- MeuFiltro.spf              (o arquivo de projeto)
|-- media_movel/               (um processador)
|   |-- Software/media_movel.cmm
|   |-- Hardware/              (gerados: media_movel.v, *.mif)
|   `-- Simulation/            (gerado: media_movel_tb.v)
|-- TopLevel/                  (Verilog de integracao, opcional)
`-- ...
\end{lstlisting}

Dois papéis especiais são atribuídos por você, pelo menu de contexto da árvore:

\begin{itemize}
  \item \textbf{Top-level} --- o módulo Verilog que integra o(s) processador(es); é o topo da síntese e o alvo do PRISM. Menu de contexto: \menu{Definir como Top Level} / \menu{Remover Top Level}.
  \item \textbf{Testbench} --- o banco de testes da simulação (um \file{_tb.v} Verilog ou um \file{.py} cocotb). Menu de contexto: \menu{Marcar como Testbench} / \menu{Desmarcar como Testbench}.
\end{itemize}

A barra de status lembra você do que falta: \enquote{Sem top-level}, \enquote{Sem testbench}. Para um projeto de um único processador, o \emph{testbench} gerado automaticamente pela compilação (\file{<proc>_tb.v}) costuma bastar --- os papéis manuais importam quando você escreve integração própria.

\section{As três visões da árvore}
\label{sec:tres-arvores}

O botão de alternância no cabeçalho do painel circula entre as visões:

\subsection{Visão Arquivos}
A visão padrão: um \emph{picker} agrupado por processador, mostrando os fontes \cmm{}, os Verilog classificados automaticamente (sintetizável vs.\ \emph{testbench}) e os artefatos relevantes. É a visão de trabalho do dia a dia.

\subsection{Visão Hierarquia}
Depois de uma síntese (por exemplo, ao abrir o PRISM), esta visão mostra a \textbf{árvore de instâncias de módulos} do projeto, como o Yosys a enxerga: o top-level no topo, cada instância aninhada abaixo. Só aparece quando há dados de hierarquia.

\subsection{Visão Pastas}
Um explorador de arquivos completo, no estilo do VS Code, com CRUD integral pelo menu de contexto:

\begin{itemize}
  \item \menu{Novo Arquivo\ldots} / \menu{Nova Pasta\ldots} / \menu{Novo arquivo Verilog (.v)} / \menu{Novo testbench cocotb (.py)} / \menu{Novo .gitignore};
  \item \menu{Recortar}, \menu{Copiar}, \menu{Colar}, \menu{Renomear\ldots}, \menu{Excluir} (com escolha entre \menu{Mover para a Lixeira} --- restaurável --- e \menu{Excluir Permanentemente});
  \item \menu{Copiar Caminho} / \menu{Copiar Caminho Relativo};
  \item \menu{Abrir no Terminal Integrado} (abre o TCMD já na pasta) e \menu{Revelar no Explorador de Arquivos};
  \item \menu{Atualizar} e \menu{Recolher Tudo}.
\end{itemize}

Nomes de arquivo são validados como o Windows exige: sem \texttt{< > : " | ? *}, sem nomes reservados do sistema, sem terminar em ponto ou espaço. Em conflitos de colagem, a AURORA oferece \menu{Substituir} ou \menu{Manter Ambos}.

\figurapendente{As três visões da árvore, lado a lado}{Três prints da mesma região: visão Arquivos, visão Hierarquia (após abrir o PRISM) e visão Pastas com o menu de contexto aberto.}

\begin{nota}
\textbf{Arrastar e soltar:} a árvore aceita a importação de arquivos \file{.v}, \file{.sv}, \file{.vh} e \emph{testbenches} cocotb \file{.py} por arrastar-e-soltar. Outros tipos são recusados com uma dica (arquivos \file{.gtkw}, por exemplo, entram pelo seletor próprio da toolbar).
\end{nota}

\section{Manutenção do projeto}

\begin{itemize}
  \item \textbf{Backup}: o botão \botao{Backup do projeto} no cabeçalho da árvore gera um \file{.zip} do projeto inteiro, sob demanda.
  \item \textbf{Arquivos faltantes}: se arquivos registrados no \file{.spf} sumirem do disco, um aviso lista quantos faltam e oferece dispensá-los (removendo as referências).
  \item \textbf{Fechar projeto}: o botão vermelho no cabeçalho da árvore, com confirmação (\enquote{Tem certeza que deseja fechar o projeto atual?}).
  \item \textbf{Renomear projeto}: disponível também via Aurora Intelligence; a operação reescreve o \file{.spf} e a pasta.
\end{itemize}
